\documentclass[11pt]{article}
\usepackage{mathblatt}


\begin{document}
\pruefheftstil{Prozent \textperiodcentered{} Portion 1}{Prüfungsheft P10}
\noindent{\Large\bfseries Prozent \textperiodcentered{} Portion 1}\hfill{\small\color{mbgrau}Prüfungsheft P10}\par\medskip
\pfrueckblick
\begin{pfnr}{}{1.}{}
In einer Tüte sind $10$ Bonbons. $6$ davon sind rot. Schreibe den Anteil der roten Bonbons als Bruch und kürze ihn.
\par\noindent \leerfeld \par
\rechenplatz{2}
\fusshilfe{1: $\tfrac{3}{5}$}
\end{pfnr}
\begin{pfnr}{}{2.}{}
Schreibe $\tfrac{1}{2}$ als Dezimalzahl.
\par\noindent \leerfeld \par
\fusshilfe{2: $0{,}5$}
\end{pfnr}
\begin{pfnr}{}{3.}{}
Berechne. $\tfrac{600}{100}$
\par\noindent \leerfeld \par
\fusshilfe{3: $6$}
\end{pfnr}
\begin{pfnr}{}{4.}{}
$2$ Kinokarten kosten $18\,€$. Berechne mit der Tabelle, wie viel $5$ Kinokarten kosten.
\par\smallskip\noindent \begin{dreisatz}{Karten}{Preis}\dsz{2}{18\,€}\dsp{\phantom{:0}}{\phantom{:0}}\dsz{1}{\dsleer}\dsp{\phantom{:0}}{\phantom{:0}}\dsz{5}{\dsleer}\end{dreisatz}\par
\par\noindent \leerfeld[€] \par
\rechenplatz{3}
\fusshilfe{4: $45\,€$}
\end{pfnr}
\begin{pfnr}{}{5.}{}
Runde $4{,}73$ auf eine Stelle nach dem Komma.
\par\noindent \leerfeld \par
\rechenplatz{1}
\fusshilfe{5: $4{,}7$}
\end{pfnr}
\begin{pfnr}{}{6.}{}
Berechne $\tfrac{1}{4}$ von $20\,€$.
\par\noindent \leerfeld[€] \par
\fusshilfe{6: $5\,€$}
\end{pfnr}
\begin{pfnr}{}{7.}{}
Berechne. $2^{6}$
\par\noindent \leerfeld \par
\fusshilfe{7: 64}
\end{pfnr}
\pfstufekopf[sprozentundanteilumwa]{Prozent und Anteil umwandeln}{in 1 der letzten 5 Prüfungen}
\pfgruppe{}
\begin{pfnr}{eigene Aufgabe}{8.}{}
In einer Stadt fahren $20\,\%$ der Berufstätigen mit dem Bus zur Arbeit. Ergänze: Jeder \leerfeld. Berufstätige fährt mit dem Bus.
\rechenplatz{2}
\fusshilfe{8: jeder $5.$}
\end{pfnr}
\pfgruppe{}
\begin{pfnr}{eigene Aufgabe}{9.}{}
Schreibe den Bruch $\tfrac{11}{20}$ in Prozent.
\par\noindent \leerfeld[\%] \par
\rechenplatz{2}
\fusshilfe{9: $55\,\%$}
\end{pfnr}
\begin{pfnr}{eigene Aufgabe}{10.}{}
Der Streifen zeigt ein Ganzes. Ein Teil davon ist grau. Wie viel Prozent des ganzen Streifens sind grau?
\par\smallskip\noindent \streifenfeld{50}{$0\,\%$}{$100\,\%$}\par
\par\noindent \leerfeld[\%] \par
\rechenplatz{1}
\fusshilfe{10: $50\,\%$}
\end{pfnr}
\pfgruppe{}
\begin{pfnr}{P10 ’22}{11.}{1 BE}
Ein Fünftel aller Jugendlichen erhält kein Taschengeld. Kreuze den passenden Prozentsatz an:
\pfkreuzzeile{5 \%}\pfkreuzzeile{20 \%}\pfkreuzzeile{25 \%}\pfkreuzzeile{50 \%}
\fusshilfe{11: 20 \%}
\end{pfnr}
\begin{pfnr}{P10 ’20}{12.}{1 BE}
Ein Förster sagt: „Bei mir wachsen nur 4 \% der gepflanzten Bäume nicht an.“ Kreuze an, welche Aussage dasselbe bedeutet:
\pfkreuzzeile{Von 10 Bäumen wachsen 4 nicht an.}\pfkreuzzeile{Jeder vierte Baum wächst nicht an.}\pfkreuzzeile{Von 100 Bäumen wachsen 4 nicht an}
\fusshilfe{12: 4 von 100 Bäumen wachsen nicht an}
\end{pfnr}
\end{document}